\subsection{Removing every field while retaining bounded degree}
\label{res:zero-field}

The field terms can also be removed without hiding a high-degree
reference spin.  A single shared reference per gadget preserves the
character symmetry; bounded-degree expansion synchronizes these
references energetically.

Let $M=|\mathcal E_n|$ and choose a 16-regular bipartite Ramanujan
graph $R_n$ on $N_R$ vertices with $M\le N_R<2M$.  Such sizes can
be obtained by successive 2-lifts starting from $K_{16,16}$, using
\href{https://arxiv.org/abs/1304.4132}{Marcus--Spielman--Srivastava}.
Its second adjacency eigenvalue is at most $2\sqrt{15}$, so its edge
expansion satisfies
\[
 |\partial B|\ge(8-\sqrt{15})|B|\ge4|B|
 \quad\text{whenever }|B|\le N_R/2.
\]
Assign a distinct reference vertex $r_e$ to each test $e=(z,w)$; any
extra reference vertices carry no gadget.  With $\mathcal U$ as above,
define the temperature-independent, zero-field pairwise Hamiltonian
\[
 \begin{split}
 H^{(0)}_n(X,S,R)={}&
 -j\sum_{e=(z,w)}\sum_{u\in\mathcal U}s_{e,u}
       (u_1X_z+u_2X_w+u_3X_{z+w}-r_e)\\
 &-j\sum_{\{v,v'\}\in E(R_n)}r_vr_{v'}.
 \end{split}
 \tag{ZF}
\]
All the terms are pair couplings, and every visible, auxiliary, and
reference spin is part of the input oracle.

\begin{theorem}[Zero-field bounded-degree full-AA failure]
\label{res:zero-field-theorem}
For every fixed
\[
 \beta j>\frac{(1+6C)\log2}{64},\qquad C=2048,
\]
the Gibbs family of (ZF) violates full acceptance simulation.  One
quantum query requires $\Omega(\log N_n)$ classical queries at
constant mean error, where $N_n=n+4M+N_R\le(1+6C)n$.  The maximum
graph degree is at most $\max\{4D,20\}$ and incident absolute
strength is at most $12Dj$.  Thus both remain bounded independently
of system size, and the fields vanish identically.
\end{theorem}

\begin{proof}
First establish an energy gap.  For a reference configuration $R$,
choose its majority sign $r\in\{-1,1\}$ and let $B$ contain its
minority vertices.  Replacing all references by $r$ changes each
occupied minority vertex's four gadget couplings by at most $8j$ in
the favorable direction.  The reference ferromagnet charges
$2j|\partial B|\ge8j|B|$.  Consequently, for every fixed $X,S$,
\[
 H^{(0)}_n(X,S,R)\ge H^{(0)}_n(X,S,r\mathbf1).
 \tag{ZF-compare}
\]

For uniform references, minimizing the four auxiliary spins of a
gadget gives energy $-8j$ if $X_zX_wX_{z+w}=r$, and $-4j$ otherwise.
This follows from
\[
 \sum_{u\in\mathcal U}|u_1x+u_2y+u_3z-r|
 =\begin{cases}8,&xyz=r,\\4,&xyz=-r.\end{cases}
\]
The identity remains valid for repeated visible indices.  Thus the
global ground energy is
$E_0=-8jM-j|E(R_n)|$, attained at least at the $2n$ choices of
uniform reference sign $r$, visible input $X=r c_a$, and minimizing
auxiliary configuration.

If $X$ has distance at least $n/32$ from every signed character
$\pm c_a$, then $rX$ is at that distance from every character.
Sparse soundness supplies at least $16n$ tests with product different
from $r$.  Combining their $4j$ cost with (ZF-compare) gives
$H^{(0)}_n\ge E_0+64jn$ for every auxiliary/reference completion.
Therefore
\[
 \Pr\{\operatorname{dist}(X,\{\pm c_a\})\ge n/32\}
 \le\frac1{2n}
 \exp\bigl(((1+6C)\log2-64\beta j)n\bigr)=:\delta_n.
 \tag{ZF-mass}
\]

It remains to justify the query lower bound despite access to the
reference spins.  Summing auxiliaries using the fourfold identity
gives the exact joint marginal
\[
 \mu(X,R)\propto
 \exp\left(K(\beta j)\sum_e r_eX_zX_wX_{z+w}
       +\beta j\sum_{\{v,v'\}\in E(R_n)}r_vr_{v'}\right).
 \tag{ZF-marginal}
\]
It is invariant under $X\mapsto Xc_b$ while $R$ is left fixed.
On the good event in (ZF-mass), write uniquely
$X=(-1)^B c_A\eta$, with $\eta$ within distance $n/32$ of the
all-plus string.  Invariance of (ZF-marginal) makes $A$ uniform and
independent of $(B,\eta,R)$, conditional on this event.  The acceptance
function $P$ is unchanged by the global sign, and
$|P(X)-h(A)|\le1/8$.

Grant the entire reference vector, $B$, and $\eta$ to the classical
simulator, and represent the conditionally independent auxiliary
spins by independent seeds.  Fixing these seeds preserves uniformity
of $A$.  A visible query uses one linear query to $A$; a reference
query uses none; an auxiliary query uses at most three.  The proof of
Lemma~\ref{res:orbit-lemma} consequently gives
\[
 \mathbb E|D-P|\ge
 (1-\delta_n)\left[\frac{1-2^{3t-m/2}}2-\frac18\right]
\]
for every all-spin depth-$t$ tree.  This proves the separation.

Visible degrees and strengths have the bounds from the fourfold
gadget.  Auxiliaries now have at most three visible neighbors and one
reference neighbor, with incident strength at most $4j$.  Reference
vertices have sixteen expander neighbors and at most four auxiliary
neighbors, giving degree and incident strength at most $20$ and
$20j$.  These are bounded by the theorem's stated maxima.
\end{proof}

The shared reference is essential.  Giving different auxiliary spins
unrelated reference vertices would let the reference graph break the
character-permutation symmetry used in (ZF-marginal).  No such
symmetry-breaking replacement is used here.
